\section{Whole quasicoherent sections localize at the original parameter}

\begin{lemma}[Every original affine tuple restriction is parameter localization]
\label{lem:scheme-module-affine-tuple-localization}
\leanok
\lean{EnvelopeProperty.SchemeModulePrincipalLocalization.factorRestriction_isLocalizedModule}
Let $X$ be a separated scheme, $\pi:X\to B$ any morphism,
$M$ a quasicoherent module on $X$, $r\in\Gamma(B,\mathcal O_B)$,
and $(U_i)$ a family of affine opens. For every nonempty finite tuple,
the original restriction from its intersection $C$ to
$C\cap\pi^{-1}D(r)$ is localization at the same original $r$,
with its original $\Gamma(B,\mathcal O_B)$-module actions.
\uses{sc0-affine-global-parameter:actual-restriction-localization}
\end{lemma}
\begin{proof}
Separatedness makes $C$ affine: finite nonempty intersections
of affine opens are affine because the diagonal is affine.
Apply the affine-open localization theorem to this very $C$
and the same original coefficient map. The tuple source at
$X$ equals $C$ and its target at $\pi^{-1}D(r)$ equals
$C\cap\pi^{-1}D(r)$. These equalities identify the actual
restriction maps and coefficient actions, yielding the claim.
\end{proof}

\begin{theorem}[Whole original section restriction is localization]
\label{thm:scheme-module-whole-principal-localization}
\leanok
\lean{EnvelopeProperty.SchemeModulePrincipalLocalization.restriction_isLocalizedModule}
If $X$ is quasicompact and separated, for every such $\pi,M,r$
the entire original restriction
\[
 \Gamma(X,M)\longrightarrow\Gamma(\pi^{-1}D(r),M)
\]
is localization at the same $r$ as a map of original
$\Gamma(B,\mathcal O_B)$-modules. Empty schemes and zero
parameters are included.
\uses{lem:scheme-module-affine-tuple-localization,prop:sheaf-finite-cech-term-localization,prop:scheme-module-cech-localization-criterion}
\end{theorem}
\begin{proof}
Choose a finite affine open cover of $X$ using quasicompactness.
In both degrees zero and one, every tuple intersection
restriction localizes by the preceding lemma. The entire
Cech term restriction therefore localizes by the finite
product theorem, using common powers of the same $r$.
The original section modules are the actual Cech kernels,
and these identifications commute with the original open
restrictions. The two-term kernel-localization theorem
therefore proves that the entire original section restriction
is localization. A finite empty cover is permitted throughout,
so no nonemptiness hypothesis has been used.
\end{proof}
